\section{核心代码}

\begin{sloppypar}%
\textbf{支撑材料文件列表}：全部代码位于 \texttt{solve/code/}，由
\texttt{common.py}（路径与常量）、\texttt{p1\_quality.py}（问题一评分与
冲突消解）、\texttt{p1\_mixture.py}（问题一配比建模）、
\texttt{p2\_scaling.py}（问题二标度律）、\texttt{p3\_optimization.py}
（问题三优化）、\texttt{p4\_evolution.py}（问题四分解预测）六个脚本与
\texttt{plotstyle.py}（绘图样式）构成，
运行环境为 Python 3.14 + numpy 2.4 + pandas 3.0 + scipy 1.17 +
scikit-learn 1.8 + matplotlib 3.10；结果写入 \texttt{solve/results/}，
插图写入 \texttt{solve/figures/}（正文全部图件）；每个实验的复算脚本位于
\texttt{solve/experiments/v*\_*.py}，对应的数值输出为
\texttt{solve/results/v*\_*.json/.csv}。以下列出五个主要脚本的核心代码
片段（\texttt{common.py} 为路径与常量工具脚本、\texttt{plotstyle.py} 为
绘图样式脚本，二者随提交文件一并提供；全部注释均为中文）。
\end{sloppypar}

\subsection{问题一：质量评分与冲突消解（\texttt{p1\_quality.py}）}

\begin{lstlisting}[language=Python, basicstyle=\ttfamily\small]
def entropy_weight(X):
    """信息熵权重: H_j 越小(分布越集中)权重越大"""
    X = np.clip(X, 1e-12, 1.0)
    P = X / X.sum(axis=0, keepdims=True)
    H = -(P * np.log(P)).sum(axis=0) / np.log(len(X))
    w = (1 - H) / np.maximum((1 - H).sum(), 1e-12)
    return w / w.sum()

def critic_weight(X):
    """CRITIC 权重: 标准差大且独立性高的指标权重大"""
    n = X.shape[1]
    sigma = X.std(axis=0, ddof=1)
    R = np.corrcoef(X.T)
    indep = np.array([1 - np.abs(R[j]).sum() / (n - 1) for j in range(n)])
    w = sigma * indep
    return w / w.sum()

def topsis(X, w):
    V = X * w
    Vp, Vn = V.max(axis=0), V.min(axis=0)
    dP = ((V - Vp) ** 2).sum(axis=1) ** 0.5
    dN = ((V - Vn) ** 2).sum(axis=1) ** 0.5
    return dN / (dP + dN)            # 贴近度 => 质量分

def conflict_index(frame):
    """双族冲突指数 = |内容质量族均值 - 格式清洁族均值|"""
    Qc = frame[FAMILY_C].mean(axis=1)
    Qf = frame[FAMILY_F].mean(axis=1)
    return (Qc - Qf).abs()

def robust_resolve(X, w, k=6):
    """对称截尾融合: 同删最低 k/2 与最高 k/2 个指标后重归一化"""
    order = np.argsort(X, axis=1)          # 每样本科升序排列
    keep = order[:, k//2 : X.shape[1]-k//2]
    sel = np.zeros_like(X, dtype=bool)
    np.put_along_axis(sel, keep, True, axis=1)
    W = np.where(sel, w, 0.0)
    return (X * W).sum(axis=1) / np.maximum(W.sum(axis=1), 1e-12)
\end{lstlisting}

\subsection{问题一：跨尺度系数收缩（\texttt{p1\_mixture.py}）}

\begin{lstlisting}[language=Python, basicstyle=\ttfamily\small]
def fit_eval(Xtr, Ytr, Xte, Yte, alpha=1.0, scale_correct=True):
    """岭回归配比模型; scale_correct 为尺度校正 R2 (去均值)"""
    model = Ridge(alpha=alpha)
    model.fit(Xtr, Ytr)
    pred = model.predict(Xte)
    if scale_correct:
        resid = (Yte - pred) - (Yte - pred).mean()   # 去整体偏差
        r2 = 1 - (resid ** 2).sum() / ((Yte - Yte.mean()) ** 2).sum()
    else:
        r2 = r2_score(Yte, pred)
    return model, r2

# 跨尺度收缩: ||beta(N)|| = c * N^{-kappa}
ns = np.array([0.001, 0.06, 1.0])          # 1M/60M/1B 规模参数量(B)
norms = np.array([norm_beta_1m, norm_beta_60m, norm_beta_1b])
kappa = -np.polyfit(np.log(ns), np.log(norms), 1)[0]
\end{lstlisting}

\subsection{问题二：广义标度律拟合（\texttt{p2\_scaling.py}）}

\begin{lstlisting}[language=Python, basicstyle=\ttfamily\small]
def fit_generalized(N, D, Q, L, form="interaction_N"):
    """最小二乘拟合 L = E + A N^-a + B D^-b + C(1-Q)^g N^-h
    (交互-N 为主链路形式: 质量缺口随参数量衰减; 附录 B 亦拟合
    additive/interaction-D/乘性等对照形式, 见 v50 BIC 表)"""
    def resid(p):
        E, A, a, B, b, C, g, h = p
        return (E + A * N**(-a) + B * D**(-b)
                + C * np.maximum(1 - Q, 0)**g * N**(-h) - L)
    p0 = np.array([L.min()*0.7, 3, 0.3, 3, 0.3, 1, 1.5, 0.3])
    lb = np.array([0, 1e-6, 1e-4, 1e-6, 1e-4, 1e-6, 1e-4, 1e-4])
    ub = np.array([L.min()*1.2, 1e3, 5, 1e3, 5, 1e3, 10, 5])
    return least_squares(resid, p0, bounds=(lb, ub), max_nfev=60000)

def r2_offset(y, yhat):
    mu = (y - yhat).mean()                 # 家族偏移修正 R2
    return 1 - ((y - yhat - mu)**2).sum() / ((y - y.mean())**2).sum()
\end{lstlisting}

\subsection{问题三：算力约束优化（\texttt{p3\_optimization.py}）}

\begin{lstlisting}[language=Python, basicstyle=\ttfamily\small]
def solve_opt(C, form, L_ctx, Q0v=0.4):
    """SLSQP 多起点求解 min L(N,D,Q) s.t. 成本<=C"""
    def obj(z):
        lnN, lnD, Q = z
        Q = np.clip(Q, Q0v, 1.0)
        return loss_generalized(np.exp(lnN), np.exp(lnD), Q)

    def cons(z):
        lnN, lnD, Q = z
        N, D = np.exp(lnN), np.exp(lnD)
        ct = 6e18 * N * D                          # C_train
        cq = D*1e9 * max(g_cost(Q, form) - g_cost(Q0v, form), 0.0)
        ca = ETA * (N*1e9) * (D*1e9) * L_ctx       # C_attn
        return C - (ct + cq + ca)

    best = None
    for seed in range(24):
        rng = np.random.default_rng(seed)
        z0 = [rng.uniform(np.log(0.01), np.log(5000)),
              rng.uniform(np.log(1), np.log(5000)),
              rng.uniform(Q0v, 0.98)]
        sol = minimize(obj, z0, method="SLSQP",
                       constraints={"type": "ineq", "fun": cons},
                       bounds=[(np.log(0.005), np.log(1e5)),
                               (np.log(0.2), np.log(1e5)), (Q0v, 1.0)],
                       options={"maxiter": 600, "ftol": 1e-12})
        if sol.success and (best is None or sol.fun < best[0]):
            best = (sol.fun, sol.x)
    return best
\end{lstlisting}

\subsection{问题四：前沿分解与预测（\texttt{p4\_evolution.py}）}

\begin{lstlisting}[language=Python, basicstyle=\ttfamily\small]
def quantile_fit(X, y, tau=0.9):
    """分位数回归主实现: statsmodels QuantReg(线性规划/单纯形),
    与 scipy.optimize.linprog(LP)、sklearn QuantileRegressor、
    L-BFGS-B、Adam 四种独立实现交叉复算 (v74 五框架, 系数一致)"""
    from statsmodels.regression.quantile_regression import QuantReg
    res = QuantReg(y, X).fit(q=tau)
    return res.params

# 增长核算: 规模贡献 = b_N * g_N, 非规模 = b_T
beta = quantile_fit(np.column_stack([np.ones(n), lnN, t]), lnS)
bN, bT = beta[1], beta[2]
gN = np.polyfit(years, np.log(params90), 1)[0]     # C4 90% 分位参数
scale_share = bN * gN / (bN * gN + bT)

# 自助法 90% 预测区间
paths = [[c0 + bN*(lnN_now + gNsc*h/12) + bT*(t_h-2022)
          + np.random.normal(0, 0.12) for _ in range(800)]
         for h in horizons]
\end{lstlisting}